\documentclass[12pt, a4paper]{article}

% Required Packages
\usepackage{amsmath}
\usepackage{amssymb}
\usepackage{kotex}
\usepackage{geometry}
\geometry{a4paper, top=1.5cm, left=2.5cm, right=2.5cm, bottom=2.5cm}

% Title
\title{linear algebra Quiz 1}
\author{5문제, 30분}
\date{ 반번호이름 : \underline{\hspace{5cm}} }

\begin{document}
\vspace{-2cm}
\maketitle 
\begin{enumerate}
    \item Prove the theorem. [10points]\\
    For an $n\times n$ matrix $A$, the following are equivalent:\\
    (a) $A$ is invertible.\\
    (b) $A\mathbf{x}=\mathbf{0}$ has only the trivial solution.\\
    (c) The reduced row echelon form of $A$ is $I_n$.\\
    (d) $A$ is a product of elementary matrices.\\
    (e) $A\mathbf{x}=\mathbf{b}$ is consistent for every $n\times1$ matrix $\mathbf{b}$.\\
    (f) $A\mathbf{x}=\mathbf{b}$ has exactly one solution for every $n\times1$ matrix $\mathbf{b}$. \\ \\ \\ \\ \\ \\ \\ \\ \\ \\ \\ \\ \\ \\ \\ \\ \\ \\ \\ \\ \\ \\ \\ \\ \\ \\ \\ \\ \\ \\ \\ \\ \\
    
    \item Prove the theorem. [10points]\\
    A system of linear equations has zero, one, or infinitely many solutions. \\ \\ \\ \\ \\ \\ \\ \\ \\ \\ \\ \\ \\ \\ \\ \\ \\ \\ \\ 
          
    \item Prove the theorem. [10points]\\
    Every linear transformation from $R^n$ to $R^m$ is a matrix transformation, and conversely every matrix transformation is a linear transformation.  \\ \\ \\ \\ \\ \\ \\ \\ \\ \\ \\ \\ \\ \\ \\ \\ \\ \\ \\ \\ \\ \\ \\ \\   
    
    \item Write the adjoint matrix of A and prove the theorem. [10points]\\
    If $A$ is an invertible $n\times n$ matrix, then
$A^{-1}=\frac{1}{\det(A)}\operatorname{adj}(A)$ \\ \\ \\ \\ \\ \\ \\ \\ \\ \\ \\ \\ \\ \\

    \item Prove the theorem. [10points] \\
    If $A\mathbf{x}=\mathbf{b}$ is a system of $n$ equations in $n$ unknowns with $\det(A)\neq 0$,
the unique solution is
\[
x_1=\frac{\det(A_1)}{\det(A)},\quad
x_2=\frac{\det(A_2)}{\det(A)},\quad\ldots,\quad
x_n=\frac{\det(A_n)}{\det(A)}
\]
where $A_j$ is the matrix obtained by replacing the $j$th column of $A$ by
$\mathbf{b}=\begin{bmatrix}b_1\\b_2\\\vdots\\b_n\end{bmatrix}$. \\ \\ \\ \\ \\ \\ \\ \\


\end{enumerate}
\end{document}